\documentclass[12pt,a4paper]{article}

\usepackage[utf8]{inputenc}
\usepackage[T1]{fontenc}
\usepackage[provide=*,spanish]{babel}
\usepackage{mathpazo}
\usepackage{microtype}
\usepackage{amsmath,amssymb,amsthm}
\usepackage{booktabs}
\usepackage{tabularx}
\usepackage{geometry}
\usepackage{enumitem}
\usepackage[hidelinks]{hyperref}

\geometry{margin=2.5cm}

\title{Ecuaciones Diferenciales Ordinarias\\Material suplementario de la Clase 2}
\author{Benjamín Marcolongo}
\date{}

\begin{document}

	\maketitle

	\section*{Introducción}

	Este material complementa la Clase 2 y reúne cuatro contenidos: ecuaciones autónomas, estabilidad, integración directa y separación de variables. Los dos primeros profundizan el análisis cualitativo de las soluciones; los dos últimos presentan métodos analíticos para obtener familias de soluciones.

	Las nociones de ecuación de primer orden, problema de valor inicial, existencia, unicidad, solución implícita y campo de direcciones se desarrollan en el material principal de la Clase 2.

	\clearpage
	\section{Ecuaciones autónomas}

	Una ecuación diferencial de primer orden es \textbf{autónoma} si la variable independiente no aparece explícitamente:

	\[
	\dot{y}=f(y).
	\]
Si la variable independiente representa el tiempo, la tasa de cambio depende solamente del estado actual del sistema y no del instante en el que se lo observa.

	\subsection{Puntos críticos y soluciones de equilibrio}

	Un valor \(y_\ast\) es un \textbf{punto crítico} o \textbf{punto de equilibrio} si

	\[
	f(y_\ast)=0.
	\]
En ese caso, la función constante

	\[
	y(x)=y_\ast
	\]
es una solución de la ecuación diferencial.

	Los puntos críticos dividen la recta de estados en intervalos. Dentro de cada intervalo, el signo de \(f(y)\) determina el sentido de evolución:

	\[
	f(y)>0 \quad\Longrightarrow\quad \dot{y}>0,
	\]
por lo que \(y\) crece, mientras que

	\[
	f(y)<0 \quad\Longrightarrow\quad \dot{y}<0,
	\]
por lo que \(y\) decrece. Esta información se representa mediante una \textbf{recta de fase}.

	\clearpage
	\section{Estabilidad}

	Un equilibrio es:

	\begin{itemize}[leftmargin=*]
		\item \textbf{asintóticamente estable} si las soluciones que comienzan suficientemente cerca permanecen cerca y tienden hacia él;
		\item \textbf{inestable} si existen soluciones inicialmente cercanas que se alejan;
		\item \textbf{semiestable} si atrae soluciones de un lado y las repele del otro.
	\end{itemize}

	En una recta de fase, un equilibrio es asintóticamente estable si las flechas apuntan hacia él desde ambos lados, e inestable si apuntan hacia afuera.

	\subsection{Ejemplo: una ecuación cúbica}

	Consideremos la ecuación autónoma

	\[
	\dot{y}=y-y^3=y(1-y)(1+y).
	\]
Los puntos críticos satisfacen

	\[
	y(1-y)(1+y)=0,
	\]
por lo que

	\[
	y_\ast=-1,\qquad y_\ast=0,\qquad y_\ast=1.
	\]
El signo de \(f(y)=y-y^3\) en los intervalos determinados por esos puntos es

	\begin{center}
		\begin{tabular}{ccccc}
			\toprule
			Intervalo & \((-\infty,-1)\) & \((-1,0)\) & \((0,1)\) & \((1,\infty)\) \\
			\midrule
			Signo de \(f(y)\) & \(+\) & \(-\) & \(+\) & \(-\) \\
			Evolución & crece & decrece & crece & decrece \\
			\bottomrule
		\end{tabular}
	\end{center}

	Por lo tanto, \(y=-1\) e \(y=1\) son equilibrios asintóticamente estables, mientras que \(y=0\) es inestable. Esta conclusión se obtiene sin resolver explícitamente la ecuación.

	\subsection{Ejemplo físico: velocidad terminal}

	Tomemos como positiva la dirección vertical hacia abajo. Para un cuerpo que cae con una resistencia del aire proporcional a la velocidad,

	\[
	m\dot{v}=mg-kv,
	\]
con \(m>0\) y \(k>0\). En forma normal,

	\[
	\dot{v}=g-\frac{k}{m}v.
	\]
El único punto crítico es

	\[
	v_T=\frac{mg}{k}.
	\]
Si \(v<v_T\), entonces \(\dot{v}>0\) y la velocidad aumenta. Si \(v>v_T\), entonces \(\dot{v}<0\) y la velocidad disminuye. En ambos casos la evolución apunta hacia \(v_T\); por eso la velocidad terminal es un equilibrio asintóticamente estable.

	\begin{quote}
		\textbf{Relación con la Guía 1.} Los problemas 15 y 16 se resuelven cualitativamente identificando puntos críticos y construyendo la recta de fase.
	\end{quote}

	Hasta aquí no fue necesario resolver explícitamente las ecuaciones. El signo de \(f(y)\), los puntos críticos y la recta de fase permitieron describir el comportamiento de las soluciones. A partir de la sección siguiente comenzaremos a buscar fórmulas mediante un método analítico.

	\clearpage
	\section{Resolución por integración directa}

	El caso más sencillo de una EDO de primer orden es

	\[
	\dot{y}=g(x),
	\]
donde la derivada depende únicamente de la variable independiente. Integrando respecto de \(x\), se obtiene

	\[
	y(x)=\int g(x)\,dx+C.
	\]
La constante \(C\) representa una familia uniparamétrica de soluciones. Una condición inicial permite seleccionar un único valor de \(C\).

	\subsection{Ejemplo: un problema de valor inicial}

	Consideremos

	\[
	\begin{cases}
		\dot{y}=x^2+1,\\
		y(0)=1.
	\end{cases}
	\]
	Integrando,\nopagebreak[4]

	\[
	y(x)=\frac{x^3}{3}+x+C.
	\]
La condición inicial implica

	\[
	1=y(0)=C.
	\]
Por lo tanto,

	\[
	y(x)=\frac{x^3}{3}+x+1.
	\]

	\clearpage
	\section{Ecuaciones de variables separables}

	Una ecuación de primer orden es \textbf{separable} si puede escribirse como

	\[
	\dot{y}=g(x)h(y).
	\]
Antes de dividir por \(h(y)\), deben buscarse sus ceros. Si

	\[
	h(y_\ast)=0,
	\]
entonces \(y(x)=y_\ast\) es una solución de equilibrio. Dividir por \(h(y)\) puede hacer que estas soluciones se pierdan.

	En un intervalo donde \(h(y)\neq 0\), podemos separar las variables:

	\[
	\frac{1}{h(y)}\,dy=g(x)\,dx.
	\]
Integrando ambos miembros,

	\[
	\int\frac{1}{h(y)}\,dy
	=
	\int g(x)\,dx+C.
	\]
	El resultado puede definir a $y$ de manera explícita o dejar una relación implícita, en el sentido desarrollado en la sección sobre soluciones explícitas e implícitas. Despejar $y$ es útil cuando puede hacerse, pero no es necesario para caracterizar las curvas solución.

	\subsection{Procedimiento}

	Para resolver una ecuación separable conviene seguir este orden:

	\begin{enumerate}[label=\arabic*.]
		\item escribirla en la forma \(\dot{y}=g(x)h(y)\);
		\item encontrar las soluciones de equilibrio resolviendo \(h(y)=0\);
		\item separar las variables en los intervalos donde \(h(y)\neq 0\);
		\item integrar ambos miembros;
		\item despejar \(y\), si es posible;
		\item imponer la condición inicial y determinar el intervalo de definición.
	\end{enumerate}

	\subsection{Ejemplo: crecimiento exponencial}

	La ecuación

	\[
	\dot{y}=3y
	\]
es separable. La solución de equilibrio es \(y=0\). Para \(y\neq 0\),

	\[
	\frac{1}{y}\,dy=3\,dx.
	\]
	\newpage
	Integrando,

	\[
	\ln|y|=3x+C.
	\]
Al exponenciar y absorber el signo en una nueva constante,

	\[
	y=Ce^{3x},
	\]
con \(C\in\mathbb{R}\). En este caso, la elección \(C=0\) recupera la solución de equilibrio.

	\subsection{Ejemplo: un PVI separable}

	Consideremos

	\[
	\begin{cases}
		\dot{y}=3y+1,\\
		y(0)=1.
	\end{cases}
	\]
La ecuación puede escribirse como

	\[
	\frac{dy}{3y+1}=dx.
	\]
Integrando,

	\[
	\frac{1}{3}\ln|3y+1|=x+C.
	\]
Equivalentemente,

	\[
	3y+1=Ke^{3x}.
	\]
La condición \(y(0)=1\) da \(K=4\). Por lo tanto,

	\[
	y(x)=\frac{4e^{3x}-1}{3}.
	\]
	\subsection{Ejemplo: crecimiento logístico}

	El modelo logístico normalizado está dado por

	\[
	\dot{y}=y(1-y).
	\]
Antes de separar variables identificamos las soluciones de equilibrio

	\[
	y(t)=0
	\qquad\text{y}\qquad
	y(t)=1.
	\]
Para una solución no constante con \(y\neq 0\) y \(y\neq 1\),

	\[
	\frac{dy}{y(1-y)}=dt.
	\]
	\newpage
	Como

	\[
	\frac{1}{y(1-y)}=\frac{1}{y}+\frac{1}{1-y},
	\]
	obtenemos

	\[
	\ln|y|-\ln|1-y|=t+C.
	\]
Luego,

	\[
	\ln\left|\frac{y}{1-y}\right|=t+C,
	\]
y la familia de soluciones no constantes puede escribirse como

	\[
	y(t)=\frac{1}{1+Ce^{-t}},
	\]
en cualquier intervalo donde el denominador no se anule. La solución de equilibrio \(y=1\) corresponde a \(C=0\), mientras que \(y=0\) debe agregarse por separado.

	El análisis cualitativo permite interpretar esta fórmula. Si \(0<y<1\), entonces \(\dot{y}>0\) y la población crece hacia \(y=1\). Si \(y>1\), entonces \(\dot{y}<0\) y decrece hacia el mismo equilibrio. Por lo tanto, \(y=1\) es asintóticamente estable y \(y=0\) es inestable.

	\begin{quote}
		\textbf{Relación con la Guía 1.} Los problemas 18--20 practican integración directa, separación de variables y PVI. El problema 25 retoma estos métodos en modelos de crecimiento.
	\end{quote}

	\clearpage
	\section*{Articulación con la Clase 2}

	Las ecuaciones autónomas y la estabilidad prolongan el análisis cualitativo iniciado con los campos de direcciones en la Clase 2. La integración directa y la separación de variables introducen luego métodos analíticos para obtener familias de soluciones explícitas o implícitas.

\end{document}